% =====================================================================
% Clase -- Álgebra 9° -- Trimestre I -- Semana 03
% Sesión 005 (mar 15 sep., 13:00, 40 min) · 006 (vie 18 sep., 6:55, 55 min)
% Tema: T1 Números reales
% Hilo conductor del trimestre I: «Al-Juarismi y el nacimiento del álgebra»
% Tema 1 de la guía: Números reales: aproximación, error y radicales (tema:numerosreales)
% Material del docente (no se entrega a los estudiantes).
% ---------------------------------------------------------------------
% El plan anual quedó aprobado el 2026-09-18 y las filas 005 y 006 de
% curriculo/programacion.csv están llenas desde el 2026-09-17: el tema, el subtema, el DBA
% y las marcas de taller y tarea de abajo salen de ahí.
% =====================================================================
\documentclass[11pt]{article}
\makeatletter\def\input@path{{../../../../../plantillas/estilo/}}\makeatother
\usepackage{mmcantillo}
\guiadelestudiante{../../guia-didactica/guia-periodo-I-reales-y-factorizacion}

\tipodocumento{Clase}
\asignatura{Álgebra}
\titulo{Semana 3: lo exacto y lo aproximado}
\grado{9°}
\periodo{I}

% DBA: PROPUESTO, sin aprobar — matematicas grado 9 · DBA 1 — «Utiliza números reales en
% sus diferentes representaciones y en diversos contextos.» Las dos sesiones trabajan la
% representación en la recta (005) y la distinción entre valor exacto y aproximación, con
% su error (006).

\begin{document}

\encabezado*

\begin{center}
{\renewcommand{\arraystretch}{1.3}
\begin{tabularx}{\linewidth}{@{}l l l X l@{}}
  \toprule
  \textbf{Sesión} & \textbf{Fecha} & \textbf{Min} & \textbf{Subtema} & \textbf{Entregables}\\
  \midrule
  005 & mar 15 sep., 13:00 & 40 & Racionales e irracionales en la recta real & ---\\
  006 & vie 18 sep., 6:55 & 55 & Aproximación y error: exactitud frente a aproximación
      ($\sqrt{2}$, $\pi$) & Taller (en clase) y tarea (para el mar 22 sep.)\\
  \bottomrule
\end{tabularx}}
\end{center}

Las semanas 01 y 02 fueron diagnóstico y refuerzo de factorización. Esta semana abre el
Tema 1 de la guía y el trimestre de verdad: el conjunto de los números reales y, sobre
todo, una distinción que 9° suele tener borrosa —un número no es lo mismo que su
aproximación decimal—. Referencia: \refguia{tema:numerosreales}.

\begin{nota}
\textbf{Pendiente.} Los ejercicios de esta semana vienen del banco de Álgebra 8°
(\texttt{irracionales-8}): falta decidir cómo se reparten los irracionales entre 8° y 9°
para no repetir el trimestre entero (ver \texttt{PENDIENTES.md}).
\end{nota}

% ---------------------------------------------------------------------
\seccion{Sesión 005 — martes 15 de septiembre · 13:00 · 40 min}

\textbf{Objetivo.} Distinguir racionales de irracionales por su expresión decimal y ubicar
raíces en la recta real con una construcción exacta.

\textbf{Materiales.} Tablero; regla y compás (al menos uno por pareja); calculadora; la
guía del estudiante.

\begin{center}
{\renewcommand{\arraystretch}{1.3}
\begin{tabularx}{\linewidth}{@{}l l X@{}}
  \toprule
  \textbf{Momento} & \textbf{Min} & \textbf{Actividad}\\
  \midrule
  Enganche & 0--5 & El recuadro del hilo del Tema 1: la tablilla babilónica YBC 7289 trae
    una aproximación de $\sqrt{2}$ con un error menor a una millonésima, y aun así
    \emph{no es} $\sqrt{2}$. Pregunta: «si la calculadora muestra $\num{1,4142136}$,
    ¿eso es $\sqrt{2}$?»\\
  Decimal y fracción & 5--15 & Repaso operativo: toda fracción da un decimal exacto o
    periódico, y todo decimal exacto o periódico se escribe como fracción. Hacer en el
    tablero $0{,}\overline{45}$ con el truco de $100x - x$. Luego el contraste:
    $2{,}101001000100001\ldots$ no tiene período, así que no es fracción. Definir
    irracional.\\
  En la recta & 15--30 & Construcción con regla y compás: la diagonal de un cuadrado de
    lado $1$ mide $\sqrt{2}$; con el compás en $0$ se lleva a la recta. Con catetos $1$ y
    $2$ sale $\sqrt{5}$. Insistir en la frase de la guía: \emph{la construcción es exacta;
    la medida con la regla, no}. Cada pareja construye $\sqrt{2}$ en el cuaderno.\\
  Práctica & 30--37 & Las dos preguntas de abajo, en el cuaderno.\\
  Cierre & 37--40 & Idea clave: \emph{los irracionales no son raros ni escasos: están en
    la diagonal de cualquier cuadrado}. Anunciar el viernes: cuánto nos equivocamos al
    aproximarlos, y el taller.\\
  \bottomrule
\end{tabularx}}
\end{center}

\textbf{Práctica de la sesión y sus respuestas.}
\begin{enumerate}
  % Ejercicio: irracionales-8-015
  \item Clasifica cada número como natural, entero, racional o irracional (puede estar en
    varios conjuntos): $\sqrt{9}$, $\sqrt{7}$, $-\frac{4}{2}$, $\frac{22}{7}$, $\pi$,
    $2{,}101001000100001\ldots$
    \\\emph{Respuesta:} $\sqrt{9} = 3$: natural, entero y racional. $\sqrt{7}$:
    irracional. $-\frac{4}{2} = -2$: entero y racional. $\frac{22}{7}$: racional. $\pi$:
    irracional. $2{,}1010010001\ldots$: irracional, porque el bloque de ceros crece y
    nunca se repite un mismo período.
  % Ejercicio: irracionales-8-010
  \item Explica cómo llevar $\sqrt{2}$ y $\sqrt{5}$ a la recta numérica con regla y
    compás, usando un cuadrado de lado $1$ y un triángulo rectángulo de catetos $1$ y $2$.
    \\\emph{Respuesta:} La diagonal del cuadrado de lado $1$ mide
    $\sqrt{1^2 + 1^2} = \sqrt{2}$; con el compás en $0$ y abertura igual a la diagonal se
    marca $\sqrt{2} \approx \num{1,41}$. Con catetos $1$ y $2$ la hipotenusa mide
    $\sqrt{1^2 + 2^2} = \sqrt{5}$, que se lleva igual hasta $\sqrt{5} \approx \num{2,24}$.
    El compás conserva la longitud: por eso el punto marcado es exactamente $\sqrt{2}$.
\end{enumerate}

% ---------------------------------------------------------------------
\seccion{Sesión 006 — viernes 18 de septiembre · 6:55 · 55 min}

\textbf{Objetivo.} Calcular el error absoluto y el error relativo de una aproximación,
juzgar cuántas cifras hacen falta en un contexto y no confundir un número con su
aproximación.

\textbf{Materiales.} Tablero; calculadora; el taller y la tarea impresos; la guía.

\begin{center}
{\renewcommand{\arraystretch}{1.3}
\begin{tabularx}{\linewidth}{@{}l l X@{}}
  \toprule
  \textbf{Momento} & \textbf{Min} & \textbf{Actividad}\\
  \midrule
  Repaso & 0--6 & Recoger la construcción del martes. Preguntar: «¿cuál es el punto
    exacto y cuál el que marcamos con la regla?»\\
  Error & 6--22 & Definición de la guía: si $a$ aproxima a $x$, error absoluto
    $= |x - a|$ y error relativo $= \frac{|x - a|}{|x|}$, que se suele dar en \%.
    Ejemplo en el tablero: $\sqrt{2} \approx \num{1,41}$ da error absoluto
    $\num{0,0042}$ y relativo cercano al \qty{0,3}{\percent}. Escribir bien grande la regla de
    notación: \emph{$\sqrt{5} = \num{2,23}$ está mal; $\sqrt{5} \approx \num{2,24}$ está
    bien}.\\
  ¿Cuántas cifras? & 22--32 & El recuadro «Planos exactos, obras aproximadas»: la diagonal
    de una baldosa de $\qty{30}{cm}$ es $30\sqrt{2}$~cm en el plano y una aproximación en
    la obra. Decidir cuántas cifras es decidir qué error se tolera: una baldosa admite un
    milímetro; la pieza de un motor, mucho menos. Pedir dos ejemplos más al grupo (una
    cancha, un medicamento).\\
  Taller & 32--50 & Repartir el taller (5 puntos), individual, con calculadora. La docente
    pasa por los puestos; recoger al final.\\
  Cierre y tarea & 50--55 & Puesta en común de la pregunta 3 del taller (¿la suma de dos
    irracionales es irracional?). Asignar la tarea, para el martes 22 de septiembre.\\
  \bottomrule
\end{tabularx}}
\end{center}

\begin{nota}
\textbf{Errores típicos de 9°.} (1) Escribir $=$ donde va $\approx$; conviene descontarlo
explícitamente en el taller desde esta semana. (2) Creer que «irracional» quiere decir
«decimal largo»: $\frac{1}{7}$ es infinito y periódico, y es racional. (3) Suponer que
operar dos irracionales da siempre un irracional; la pregunta 3 del taller existe para
romper eso.
\end{nota}

\end{document}
